\documentclass[10pt,a4paper]{amsart}
\usepackage[margin=19mm]{geometry}
\usepackage{amsmath,amssymb,mathtools}
\usepackage[scr=rsfs,cal=euler]{mathalfa}
\usepackage{xfrac}
\usepackage[unicode,hidelinks,bookmarks=false]{hyperref}
\newcommand{\ie}{i.e.\ }
\newcommand{\cf}{cf.\ }
\newcommand{\resp}{respectively\ }
\newcommand{\Subsection}{\subsection}
\providecommand{\nobreakdash}{\nobreak\hbox{-}}
\setlength{\emergencystretch}{2em}
\multlinegap0pt
\usepackage{bm}
\usepackage{bbm}
\usepackage{dsfont}
\usepackage{xspace}
\usepackage{cancel}
\usepackage{mathtools}
\usepackage{amsthm}
\makeatletter
\@ifundefined{mycor}{\newtheorem{mycor}{Corollary}}{}
\@ifundefined{mythm}{\newtheorem{mythm}{Theorem}}{}
\makeatother
\newcommand{\B}{\mathcal{B}}
\newcommand{\e}{\epsilon}
\title[The cogrowth inequality from Whitehead's algorithm]{The cogrowth inequality from Whitehead's algorithm}
\author{Asif Shaikh}
\date{Source: arXiv:2407.16523v2 (CC BY 4.0)}
\begin{document}
\maketitle

\noindent\emph{Source and licence.} The cogrowth inequality from Whitehead's algorithm. Version arXiv:2407.16523v2 (CC BY 4.0), distributed under the Creative Commons Attribution 4.0 International licence, as recorded in the arXiv licence metadata for this exact version and shown on the abstract page https://arxiv.org/abs/2407.16523v2. Full rights notice: licenses/arxiv-2407.16523-CC-BY-4.0.txt.\par\smallskip

\noindent\textit{Editorial scope.} Counted unit: the Mythm whose source label is \texttt{thm:main\_lambda<lambda1} in the article (source0.tex lines 505--508, source label \texttt{thm:main\_lambda<lambda1}), delivered together with the article's own complete original proof of it (source0.tex lines 509--537, one \texttt{proof} environment of 29 lines). No mathematics is written, repaired or re-derived: statement and proof are the article's own text line for line. Required context retained uncounted: thm:B\_from\_BH (lines 376--378); cor:M\_partition (lines 449--463); cor:M1\_from\_M (lines 485--487); thm:Perron-Frobenius (lines 498--501). Every definition, display and label of those lines is retained unchanged. Omitted from this package and not redistributed: 1--375, 379--448, 464--484, 488--497, 502--504, 538--683. Nothing inside a retained excerpt is omitted or altered: every retained body is delivered exactly as the article's lines are written, so no omission receipt of any kind accompanies this package. Citations: the retained excerpts carry no citation command at all, so no bibliography entry of the article is transcribed and no reference is redistributed. Reference adaptation: none was needed and none was made; every reference inside the delivered unit resolves inside it, so no unresolved reference or citation is produced. Printed slips kept literal: no printed word, symbol, hypothesis or number of the counted statement or of its proof was altered. Layout and class: the article's own class is \texttt{article} with packages including inputenc, titlesec, babel, pst-plot, nicematrix, caption, authblk, subcaption, amsthm, amssymb, amsmath, amsfonts, calrsfs, upgreek; the wrapper is the lane's pinned \texttt{amsart} editorial wrapper at 10pt on A4, which restates the theorem-like environments the retained text needs and adds the article's own macro definitions that the retained text needs (\texttt{\textbackslash B}, \texttt{\textbackslash e}), with extra wrapper packages (bm, bbm, dsfont, xspace, cancel, mathtools). The delivered numbering is the unit's own. Assets: no retained span contains a figure environment, a graphics-inclusion command, a PDF-page inclusion, a table or a TikZ picture; no image file is copied into this package, the package declares no asset dependency and no image file is redistributed with it. Licence: version arXiv:2407.16523v2 of this article is distributed under the Creative Commons Attribution 4.0 International licence, as recorded in the arXiv licence metadata for this exact version and shown on the abstract page https://arxiv.org/abs/2407.16523v2. A full rights notice is delivered at licenses/arxiv-2407.16523-CC-BY-4.0.txt. Characters: every retained excerpt is pure ASCII, so the delivered engine, XeLaTeX with its default Latin Modern text font, reports no missing character.
\begin{mythm}\label{thm:B_from_BH}  
Let $H\leq F_m$ be a non-cyclic free factor, and suppose the Whitehead graph of $H$ contains a cut vertex. Then there is a Whitehead automorphism $\varphi$ such that the ergodic automaton $\B_{\varphi(H)}$ that recognizes $L_{\varphi(H)}$ can be obtained from $\B_H$ by collapsing certain edges of $\B_H$.
\end{mythm}

     \begin{mycor}\label{cor:M_partition} The matrix $M$ can be decomposed as follows:
            \begin{equation}
                M = \left(
            \begin{array}{cc}
             M' & U \\
             Z & O \\
            \end{array}
            \right), \label{eqn:decomposition_M}
            \end{equation}
            where $M'$, $U$, $Z$, and $O$ are sub-matrices with dimensions $|Q|\times|Q|$, $|Q|\times|S|$, $|S|\times|Q|$, and $|S|\times|S|$, respectively. This decomposition satisfies the following properties:
            \begin{enumerate}
                \item \label{cor:C}The rows of the matrix $U$ contain either all zeros or exactly one non-zero entry equal to 1.
                \item \label{cor:O}The matrix $O$ is a zero matrix.
            \end{enumerate}
    \end{mycor}

    \begin{mycor}\label{cor:M1_from_M}
    Let $M$ be the adjacency matrix of the automaton $\B_H$, and let $M_1$ be the $|Q|\times |Q|$ matrix obtained by applying the transformation procedure described above. Then $M_1$ is the adjacency matrix of the transition diagram of $\B$, which is isomorphic to $\B_{\varphi(H)}$.
    \end{mycor}

    \begin{mythm}\label{thm:Perron-Frobenius} (Perron-Frobenius). Suppose that $N$ is an irreducible, non-negative integral matrix. Then there is a unique positive eigenvector $\overrightarrow{w}$ of norm one for $N$, and its associated eigenvalue satisfies $\eta \geq 1$. If $\eta = 1$, then $N$ is a transitive permutation matrix. Moreover if $\overrightarrow{u}$ is a positive vector and $\beta > 0$ satisfies 
    $(N\overrightarrow{u})_i \leq \beta\overrightarrow{u}_i$ for each $i$ and 
    $(N\overrightarrow{u})_j < \beta\overrightarrow{u}_j$ for some $j$, then $\eta < \beta.$
    \end{mythm}

    \begin{mythm} \label{thm:main_lambda<lambda1}
        Let $\lambda$ and $\lambda_1$ be the Perron-Frobenious eigenvalue of $M$ and $M_1$, respectively, where $M$ and $M_1$ are as described in Corollary \ref{cor:M1_from_M}. 
        Then $ \lambda < \lambda_1.$
    \end{mythm}

    \begin{proof} 
    Recall that $Q_H = Q \cup S$. We use the notation $(q, x_k^{\e})$ for states in $Q$ and $(s, a^{\e})$ for states in $S$. Let $R_{(v, x_k^{\e})}$ and $C_{(v, x_k^{\e})}$ denote the rows and columns of the matrix $M$ corresponding to the state $(v, x_k^{\e}) \in Q_H$, and let $R'_{(q, x_k^{\e})}$ and $C'_{(q, x_k^{\e})}$ denote the rows and columns of the matrix $M_1$ corresponding to the state $(q, x_k^{\e}) \in Q$.

    Choose a positive vector $\overrightarrow{v}$ in $\mathbb{R}^{|Q|}$ so that $M_1\overrightarrow{v} = \lambda_1\overrightarrow{v}$, and let $\overrightarrow{u}$ be the vector in $\mathbb{R}^{|Q|+|S|}$ defined by $u_{(q,x_k^{\e})} = v_{(q,x_k^{\e})}$ for ${(q,x_k^{\e})} \in Q$. We will determine the remaining $|S|$ components of $\overrightarrow{u}$ based on the block decomposition of $M$ described in Corollary  \ref{cor:M_partition}.
    
    By Theorem \ref{thm:B_from_BH}, in the column $C_{(s,a^{\e})}$ in $M$, there are $d$ entries with the value $1$ and the rest are $0$. The rows in which these $d$ entries appear are $R_{o(e_i)}$, respectively, where $e_i, i = 1,\cdots,d$, and are the edges in the transition graph of $\B_H$ terminating at $(s,a^{\e})$ with label $a^{\e}.$ For each $o(e_i) \in Q, i = 1,\cdots,d$, we use Corollary \ref{cor:M_partition} to write 
     $$(M \overrightarrow{u})_{o(e_i)} = b_i + u_{(s,a^{\e})},$$  
     where $b_i = \displaystyle \sum_{(q,x_k^{\e'}) \in Q} m_{o(e_i)(q,x_k^{\e'})} u_{(q,x_k^{\e'})}.$ Also, for $(s,a^{\e}) \in S$, we write 
     $$(M \overrightarrow{u})_{(s,a^{\e})} = \sum_{(q,x_k^{\e'}) \in Q} m_{(s,a^{\e})(q,x_k^{\e'})} u_{(q,x_k^{\e'})} = b_{(s,a^{\e})}.$$
     
    By Corollary \ref{cor:M1_from_M}, the corresponding rows $R'_{o(e_i)}$ of $M_1$ can be obtained from $M$ by applying row addition 
    $$R_{o(e_i)} + R_{(s,a^{\e})} \rightarrow R_{o(e_i)},$$ 
    where $i = 1, \dots, d$, and removing rows $R_{(s,a^{\e})}$ and columns $C_{(s,a^{\e})}$ from the resultant, where $(s,a^{\e}) \in S$. 

    Using the above relationships between rows $R_{o(e_i)}, R_{(s,a^{\e})},$ and the column $C_{(s,a^{\e})}$ of $M$, and the corresponding rows $R'_{o(e_i)}$ of $M_1$, for $i = 1,\cdots,d$, we obtain
    $$(M \overrightarrow{u})_{o(e_i)} + (M \overrightarrow{u})_{(s,a^{\e})} - u_{(s,a^{\e})} = (M_1 \overrightarrow{v})_{o(e_i)} = \lambda_1 u_{o(e_i)}.$$

    Substituting the expressions for $(M \overrightarrow{u})_{o(e_i)}$ and $(M \overrightarrow{u})_{(s,a^{\e})}$, we obtain
    $$b_{(s,a^{\e})} = \lambda_1 u_{o(e_i)} - b_i.$$
    Since $b_i > 0$ and $\lambda_1 u_{o(e_i)} > b_i$, it follows that $0 < b_{(s,a^{\e})} < \lambda_1 u_{o(e_i)}.$
    To ensure that $M \overrightarrow{u} \leq \lambda_1 \overrightarrow{u}$, we require
    $$(M \overrightarrow{u})_{o(e_i)} = b_i + u_{(s,a^{\e})} \leq \lambda_1 u_{o(e_i)} \quad \text{and} \quad (M \overrightarrow{u})_{(s,a^{\e})} = b_{(s,a^{\e})} \leq \lambda_1 u_{(s,a^{\e})}.$$
    Thus, we need to choose $u_{(s,a^{\e})}$ such that $$\frac{b_{(s,a^{\e})}}{\lambda_1} \leq u_{(s,a^{\e})} \leq b_{(s,a^{\e})}.$$
    This interval is non-empty, since $b_{(s,a^{\e})} > 0$ and $\lambda_1 > 1$. Therefore, for each $(s,a^{\e}) \in S$, we choose 
    $u_{(s,a^{\e})} = \frac{b_{(s,a^{\e})}}{\lambda_1}$. Then, 
    $$(M \overrightarrow{u})_{(s,a^{\e})} = b_{(s,a^{\e})} = \lambda_1 u_{(s,a^{\e})}, \quad (M \overrightarrow{u})_{o(e_i)} = b_i + u_{(s,a^{\e})} < \lambda_1 u_{o(e_i)}.$$
    Hence, for this choice, we have $(M \overrightarrow{u})_{(v, x^{\e})} \leq \lambda_1 u_{(v, x^{\e})}$ for all $_{(v, x^{\e})} \in Q_H$, with strict inequality for $(v, x^{\e}) = o(e_i),$ where $ i = 1,\cdots, d$. 
    By Theorem \ref{thm:Perron-Frobenius}, we conclude that $\lambda < \lambda_1.$ This completes the proof.
\end{proof}

\end{document}
